\documentclass[9pt,twocolumn, twoside]{chaos}

\articletype{inv} % article type
\usepackage{lipsum}
\usepackage{tikz,xcolor,hyperref}
\usepackage[export]{adjustbox}

\definecolor{lime}{HTML}{A6CE39}
\DeclareRobustCommand{\orcidicon}{%
	\begin{tikzpicture}
	\draw[lime, fill=lime] (0,0) 
	circle [radius=0.16] 
	node[white] {{\fontfamily{qag}\selectfont \tiny ID}};
	\draw[white, fill=white] (-0.0625,0.095) 
	circle [radius=0.007];
	\end{tikzpicture}
	\hspace{-2mm}
}

\foreach \x in {A, ..., Z}{%
	\expandafter\xdef\csname orcid\x\endcsname{\noexpand\href{https://orcid.org/\csname orcidauthor\x\endcsname}{\noexpand\orcidicon}}
}

\newcommand{\orcidauthorA}{0000-0001-9151-3052}
\newcommand{\orcidauthorB}{0000-0001-9151-3052}

\vspace{4cm}

\title{Chaos Theory and Applications (CHTA) in Applied Sciences and Engineering} 

\author[$\ast$,1]{Author One\orcidA{}}
\author[$\alpha$,2]{Author Two\orcidB{}}
\author[$\beta$,3]{Author Three\orcidA{}}
\author[$\S$,4]{Author Four\orcidA{}}
\author[$\ast\ast$,5]{Author Five\orcidA{}}

\affil[$\ast$]{Author one affiliation}
\affil[$\alpha$]{Author two affiliation}
\affil[$\beta$]{Author three affiliation}
\affil[$\S$]{Author four affiliation}
\affil[$\ast\ast$]{Author five affiliation}

\keywords{Keyword \\ Keyword2 \\ Keyword3 \\  Keyword4 \\ Keyword5\\...}

\runningtitle{CHTA Journal Template} % For use in the footer 

\runningauthor{FirstAuthorLastname \textit{et al.}}


\begin{abstract}
	The abstract should be written for people who may not read the entire paper, so it must stand on its own. Your abstract should provide a synopsis of the entire article;
	begin with the broad context of the study, followed by specific background for the study;
	describe the purpose, methods and procedures, core findings and results, and conclusions of the study;
	emphasize new or important aspects of the research;
	engage the broad readership of CHTA and be understandable to a diverse audience (avoid using jargon);
	be a single paragraph of less than 250 words;
	contain the full name of the organism studied;
	NOT contain citations or abbreviations.
\end{abstract}


\begin{document}
	
	\maketitle
	
\setlength{\extrarowheight}{3pt}

	\thispagestyle{firststyle}
	\logomark
	\articletypemark
	\marginmark
	\firstpagefootnote
	
	\correspondingauthoraffiliation{1}{abc@alumnos.tr}
	\correspondingauthoraffiliation{2}{xx@academicos.it}
	\correspondingauthoraffiliation{3}{aaaa@academicos.fr}
	\correspondingauthoraffiliation{4}{trtrtr@ctb.upm.es}
	\correspondingauthoraffiliation{5}{aaa@academicos.udg.mx (\textbf{Corresponding author}).}
	
	\vspace{-44pt}
	
	\setcounter{page}{1}

	\section{Introduction}
	
	For the introduction, authors should be mindful of the broad readership of the journal. The introduction should set the stage for the importance of the work to a generalist reader and draw the reader in to the specific study. The scope and impact of the work should be clearly stated.
	
 
	Authors are encouraged to:
	
	\begin{itemize}
		\item cite the supporting literature completely rather than select a subset of citations;
		\item provide important background citations, including relevant review papers (to help orient the non-specialist reader);
		\item to cite similar work in other chaos theory and applications.
	\end{itemize}
	
	
	
		\section{First level section header}
	
	Use this level to group two or more closely related headings in a long article.
	
	\subsection{Second level section header}
	
	Second level section text.
	
	\subsubsection{Third level section header:}
	
	Third level section text. These headings may be numbered, but only when the numbers must be cited in the text. 
	\section{In-text Citations}
	
	Add citations using the \verb|\citep{}| command, for example \citep{Akgul2016} or for multiple citations, \citep{C.Sprott2000, rossler1976equation}.
	
	\section{Figures and Tables}
	
	Figures and Tables should be labelled and referenced in the standard way using the \verb|\label{}| and \verb|\ref{}| commands.
	
	\subsection{Sample Figure}
	
	Figure \ref{fig:spectrum} shows an example figure.
	
	\begin{figure*}[htbp]
		\renewcommand{\familydefault}{\sfdefault}\normalfont
		\centering
		\includegraphics[width=\linewidth]{ChaosV.pdf}
		\caption{Example figure}%
		\label{fig:spectrum}
	\end{figure*}
	
	
	\subsection{Sample Table}
	
	Table \ref{tab:shape-functions} shows an example table. 
	
	\begin{table*}[htbp]
		\renewcommand{\familydefault}{\sfdefault}\normalfont
		\centering
		\caption{\bf Example table}
		\begin{tableminipage}{\textwidth}
			\begin{tabularx}{\textwidth}{XXXX}
				\hline
				\header Student & Grade & Rank & Notes \\
				\hline
				Alice & 82\% & 1 & Performed very well.\\
				Bob & 65\% & 3 & Not up to his usual standard.\\
				Charlie & 73\% & 2 & A good attempt.\\
				\hline
			\end{tabularx}
			\label{tab:shape-functions}
		\end{tableminipage}
	\end{table*}
	
	\section{Sample Equation}
	
	
	\begin{equation}
	S_n = \frac{X_1 + X_2 + \cdots + X_n}{n}
	= \frac{1}{n}\sum_{i}^{n} X_i
	\label{eq:refname1}
	\end{equation}

	
	\section{Results and Discussion}
	
	The results and discussion should not be repetitive. The results section should give a factual presentation of the data and all tables and figures should be referenced; the discussion should not summarize the results but provide an interpretation of the results, and should clearly delineate between the findings of the particular study and the possible impact of those findings in a larger context. 
	
	\subsection*{Acknowledgments}
	
This work was supported by X University Scientific Research Projects Unit
under Grants 0000-00-00-001, 0000-00-00-002.
	
	\subsection*{Availability of data and material}
Not applicable.

\subsection*{Conflicts of interest}

The authors declare that there is no conflict of interest regarding the publication of this paper.
	
	\bibliography{example-bibliography}
	
\subsubsection{How to cite this article:}

Torre, J. O., Jaimes, R., Garcia, J. H., Pisarchik, A. N., and Alex, G. Paper Name. \textit {Chaos Theory and Applications}, x(x), xxx-xxx, 2024.

\subsubsection{Licensing Policy:}
The published articles in CHTA are licensed under a \href{https://creativecommons.org/licenses/by-nc/4.0/} {Creative Commons Attribution-NonCommercial 4.0 International License}.

\includegraphics[width=2.1cm,center]{indir.png}

\end{document}